% GENERATED FILE. Edit canonical lessons, then run npm run generate:books.
% canonical-source-hash: fnv1a64:f1cb20d763ff6ee9
% canonical-lessons: RU-C297-travma, RU-C297-ushib, RU-C297-soznaniye, RU-C297-postradat, RU-C297-ranit

\chapter{Injuries}
\label{ch:ru-injuries}
\hlchaptermodality{297}

\begin{chapteropening}
\textbf{By the end of this chapter:} \emph{I can talk about an injury, a bruise and losing consciousness, and say someone was hurt or wounded.}

Complete the last lesson of chapter 297: I can talk about an injury, a bruise and losing consciousness, and say someone was hurt or wounded.
\end{chapteropening}

\section[\texorpdfstring{\ru{травма} (trávma)}{травма (trávma)}]{\ru{травма} (trávma) — an injury}

\label{lesson:RU-C297-travma}



\begin{quote}
Before the new one: say the Russian for to evacuate, then the Russian for an evacuation.
\end{quote}

\subsection*{You'll want to know: \ru{травма}}
\textbf{\ru{травма}} (\emph{trávma}) — \textquotedblleft{}an injury\textquotedblright{}.

\begin{quote}
\ru{Футболист} \ru{получил} \ru{травму}.
\end{quote}

\begin{quote}
\emph{Futbolíst poluchíl trávmu.}
\end{quote}

\begin{quote}
The footballer got injured.
\end{quote}

Feminine. \emph{\ru{Травмпункт}} — an emergency clinic for injuries.

\begin{cousinweb}
\emph{\ru{травма}} goes back to Greek \emph{traûma} \textquotedblleft{}a wound\textquotedblright{}; EtymDB lists the Russian word with no parent.
\end{cousinweb}

\begin{grammarlens}[title={What you've built}]
One more word for telling, explaining and describing at B1.
\end{grammarlens}

\subsection*{Your turn}
\begin{itemize}
\raggedright
  \item \emph{Say it:} \emph{trávma}
  \item \emph{Say it:} \emph{trávma}, once more
  \item \emph{Say it:} a sentence of your own with \emph{trávma}
  \item \emph{Recall:} say the Russian for to evacuate, then the Russian for an evacuation, then say \emph{trávma} again
\end{itemize}

\begin{tcolorbox}[breakable,colback=teal!4,colframe=teal!35!black,title={Before you move on}]
What does \ru{травма} mean? (\textquotedblleft{}An injury\textquotedblright{}.) Say it once more.
\end{tcolorbox}
\section[\texorpdfstring{\ru{ушиб} (ushíb)}{ушиб (ushíb)}]{\ru{ушиб} (ushíb) — a bruise, a knock}

\label{lesson:RU-C297-ushib}



\begin{quote}
Before the new one: say the Russian for an evacuation, then the Russian for an injury.
\end{quote}

\subsection*{You'll want to know: \ru{ушиб}}
\textbf{\ru{ушиб}} (\emph{ushíb}) — \textquotedblleft{}a bruise, a knock\textquotedblright{}.

\begin{quote}
\ru{У} \ru{мальчика} \ru{сильный} \ru{ушиб} \ru{ноги}.
\end{quote}

\begin{quote}
\emph{U málchika sílnyy ushíb nogí.}
\end{quote}

\begin{quote}
The boy has a bad bruise on his leg.
\end{quote}

Masculine; genitive \emph{\ru{ушиба}}.

\begin{cousinweb}
\emph{\ru{ушиб}} is built on \emph{\ru{ушибить}} \textquotedblleft{}to knock, to bruise\textquotedblright{}, from the old verb \emph{\ru{шибать}} \textquotedblleft{}to strike\textquotedblright{}; EtymDB lists both with no parent.
\end{cousinweb}

\begin{grammarlens}[title={What you've built}]
One more word for telling, explaining and describing at B1.
\end{grammarlens}

\subsection*{Your turn}
\begin{itemize}
\raggedright
  \item \emph{Say it:} \emph{ushíb}
  \item \emph{Say it:} \emph{ushíb}, once more
  \item \emph{Say it:} a sentence of your own with \emph{ushíb}
  \item \emph{Recall:} say the Russian for an evacuation, then the Russian for an injury, then say \emph{ushíb} again
\end{itemize}

\begin{tcolorbox}[breakable,colback=teal!4,colframe=teal!35!black,title={Before you move on}]
What does \ru{ушиб} mean? (\textquotedblleft{}A bruise, a knock\textquotedblright{}.) Say it once more.
\end{tcolorbox}
\section[\texorpdfstring{\ru{сознание} (soznániye)}{сознание (soznániye)}]{\ru{сознание} (soznániye) — consciousness}

\label{lesson:RU-C297-soznaniye}



\begin{quote}
Before the new one: say the Russian for an injury, then the Russian for a bruise.
\end{quote}

\subsection*{You'll want to know: \ru{сознание}}
\textbf{\ru{сознание}} (\emph{soznániye}) — \textquotedblleft{}consciousness\textquotedblright{}.

\begin{quote}
\ru{Он} \ru{потерял} \ru{сознание} \ru{на} \ru{минуту}.
\end{quote}

\begin{quote}
\emph{On poteryál soznániye na minútu.}
\end{quote}

\begin{quote}
He lost consciousness for a minute.
\end{quote}

\emph{\ru{Потерять} \ru{сознание}} — to faint; \emph{\ru{прийти} \ru{в} \ru{сознание}} — to come round.

\begin{cousinweb}
\emph{\ru{сознание}} is built on \emph{\ru{сознать}} \textquotedblleft{}to realize\textquotedblright{}, from \emph{\ru{со}-} and \emph{\ru{знать}} \textquotedblleft{}to know\textquotedblright{}, inherited from Proto-Slavic \emph{znati}, PIE \emph{ǵneh\textsubscript{3}-} (EtymDB).
\end{cousinweb}

\begin{grammarlens}[title={What you've built}]
One more word for telling, explaining and describing at B1.
\end{grammarlens}

\subsection*{Your turn}
\begin{itemize}
\raggedright
  \item \emph{Say it:} \emph{soznániye}
  \item \emph{Say it:} \emph{soznániye}, once more
  \item \emph{Say it:} a sentence of your own with \emph{soznániye}
  \item \emph{Recall:} say the Russian for an injury, then the Russian for a bruise, then say \emph{soznániye} again
\end{itemize}

\begin{tcolorbox}[breakable,colback=teal!4,colframe=teal!35!black,title={Before you move on}]
What does \ru{сознание} mean? (\textquotedblleft{}Consciousness\textquotedblright{}.) Say it once more.
\end{tcolorbox}
\section[\texorpdfstring{\ru{пострадать} (postradát)}{пострадать (postradát)}]{\ru{пострадать} (postradát) — to be hurt, to suffer}

\label{lesson:RU-C297-postradat}



\begin{quote}
Before the new one: say the Russian for a bruise, then the Russian for consciousness.
\end{quote}

\subsection*{You'll want to know: \ru{пострадать}}
\textbf{\ru{пострадать}} (\emph{postradát}) — \textquotedblleft{}to be hurt, to suffer\textquotedblright{}.

\begin{quote}
\ru{В} \ru{аварии} \ru{никто} \ru{не} \ru{пострадал}.
\end{quote}

\begin{quote}
\emph{V avárii niktó ne postradál.}
\end{quote}

\begin{quote}
Nobody was hurt in the accident.
\end{quote}

Perfective. It takes \emph{\ru{от}} + genitive: \emph{\ru{пострадать} \ru{от} \ru{пожара}}.

\begin{cousinweb}
\emph{\ru{пострадать}} is built on \emph{\ru{по}-} and \emph{\ru{страдать}} \textquotedblleft{}to suffer\textquotedblright{}, an inherited Slavic verb; EtymDB lists \emph{\ru{страдать}} \textquotedblleft{}to suffer\textquotedblright{} with no parent.
\end{cousinweb}

\begin{grammarlens}[title={What you've built}]
One more word for telling, explaining and describing at B1.
\end{grammarlens}

\subsection*{Your turn}
\begin{itemize}
\raggedright
  \item \emph{Say it:} \emph{postradát}
  \item \emph{Say it:} \emph{postradát}, once more
  \item \emph{Say it:} a sentence of your own with \emph{postradát}
  \item \emph{Recall:} say the Russian for a bruise, then the Russian for consciousness, then say \emph{postradát} again
\end{itemize}

\begin{tcolorbox}[breakable,colback=teal!4,colframe=teal!35!black,title={Before you move on}]
What does \ru{пострадать} mean? (\textquotedblleft{}To be hurt, to suffer\textquotedblright{}.) Say it once more.
\end{tcolorbox}
\section[\texorpdfstring{\ru{ранить} (ránit)}{ранить (ránit)}]{\ru{ранить} (ránit) — to wound}

\label{lesson:RU-C297-ranit}



\begin{quote}
Before the new one: say the Russian for consciousness, then the Russian for to be hurt.
\end{quote}

\subsection*{You'll want to know: \ru{ранить}}
\textbf{\ru{ранить}} (\emph{ránit}) — \textquotedblleft{}to wound\textquotedblright{}.

\begin{quote}
\ru{Его} \ru{ранили} \ru{в} \ru{руку}.
\end{quote}

\begin{quote}
\emph{Yevó ránili v rúku.}
\end{quote}

\begin{quote}
He was wounded in the arm.
\end{quote}

One form for both aspects. \emph{\ru{Раненый}} — wounded.

\begin{cousinweb}
\emph{\ru{ранить}} is built on \emph{\ru{рана}} \textquotedblleft{}a wound\textquotedblright{}, inherited from Proto-Slavic \emph{rana}; EtymDB lists the verb with no parent.
\end{cousinweb}

\begin{grammarlens}[title={What you've built}]
One more word for telling, explaining and describing at B1.
\end{grammarlens}

\subsection*{Your turn}
\begin{itemize}
\raggedright
  \item \emph{Say it:} \emph{ránit}
  \item \emph{Say it:} \emph{ránit}, once more
  \item \emph{Say it:} a sentence of your own with \emph{ránit}
  \item \emph{Recall:} say the Russian for consciousness, then the Russian for to be hurt, then say \emph{ránit} again
\end{itemize}

\begin{tcolorbox}[breakable,colback=teal!4,colframe=teal!35!black,title={Before you move on}]
What does \ru{ранить} mean? (\textquotedblleft{}To wound\textquotedblright{}.) Say it once more.
\end{tcolorbox}
